%% region NMDAR dependence, ca8 ablation, LTM with puromycin

% region MK-801
% ** Delay and Trace Conditioning Performance with MK-801

\subsection{Associative Learning is Dependent on NMDA Receptor Function}

To probe the molecular basis of associative learning, I examined the role of NMDARs using the non-competitive antagonist MK-801. The experiment comprised two paradigms: delay conditioning, with MK-801 at 0~\si{\micro\Molar}, 10~\si{\micro\Molar}, or 100~\si{\micro\Molar}, and 3sTrace conditioning, with 0~\si{\micro\Molar} and 100~\si{\micro\Molar}. All test solutions contained 0.2\% DMSO. Control groups received explicitly unpaired CS--US presentations to dissociate associative from non-associative effects.

\paragraph{Delay and Trace Conditioning Performance with MK-801}
Blocking NMDARs with high-dose MK-801 (100~\si{\micro\Molar}) completely abolished CR acquisition in both delay and 3sTrace paradigms (Fig. X). In the Early Test phase, normalized vigor in these larvae did not differ from Pre-Train or from unpaired controls, indicating a failure to learn. In contrast, 10~\si{\micro\Molar} MK-801 produced no measurable impairment in delay conditioning (Fig. X), demonstrating a clear dose dependence of NMDAR involvement.
During delay conditioning, larvae exposed to 0~\si{\micro\Molar} or 10~\si{\micro\Molar} MK-801 displayed a consistent reduction in scaled vigor following CS onset from the second training block onward—an expression of a CR that extinguished by the second test block (Fig. X). In Early Test trials, CR duration was slightly prolonged in the 10~\si{\micro\Molar} group (lasting $\approx$~5~s after CS offset) relative to the 0~\si{\micro\Molar} group ($\approx$~2.5~s). By contrast, 100~\si{\micro\Molar} MK-801 disrupted learning entirely: scaled vigor remained at baseline and did not differ from unpaired controls across experimental phases (Figs. X--Y).
In the 3sTrace paradigm, larvae at 0~\si{\micro\Molar} successfully acquired a CR, whereas those treated with 100~\si{\micro\Molar} failed to do so (Fig. X). Early Test normalized-vigor values of the high-dose group were indistinguishable from Pre-Train and control values. This complete loss of associative suppression confirms that NMDAR signaling is required for both contiguous (delay) and temporally discontiguous (trace) learning.
Subtle trends toward reduced vigor in some high-dose delay fish suggested incomplete blockade or partial compensation, but these were inconsistent across individuals. Variability among 100~\si{\micro\Molar} unpaired controls likely reflects direct motor effects of MK-801 rather than genuine learning. Only delay-conditioning larvae were pre-incubated with MK-801 prior to training, a difference that may further enhance sensitivity to the drug in that group.
% todo These results led me to conclude that NMDA receptor function is critical for this form of associative learning in larval zebrafish, as expected from the literature (Hinz et al., 2013).

\paragraph{Motor Effects of MK-801}
Because I measured the CR as a reduction in tail movement vigor, maintaining a high baseline of spontaneous movement was essential for detecting learning. Sustaining sufficient baseline activity after the \SI{3}{\hour} retention interval posed a major challenge. To address this, I introduced a second priming phase immediately before testing, applying the same principle as the initial priming phase in Chapter 3. Based on pilot experiments, I delivered four \SI{100}{\milli\second} US pulses immediately before the test phase to boost tail activity and enable more reliable CR detection.

% endregion MK-801


% region ca8 ablation
% ** Purkinje Cell Ablation Impairs Associative Learning in Delay and Trace Conditioning
\subsection*{Purkinje Cell Ablation Impairs Associative Learning in Delay and Trace Conditioning}

To examine the role of cerebellar Purkinje cells (PCs) in associative learning, we performed chemogenetic ablation using Nifurpirinol in ca8$^+$ larvae. Six groups were compared: unpaired control ca8$^+$ and ca8$^-$; delay ca8$^+$ and ca8$^-$; trace ca8$^+$ and ca8$^-$.

At the population level, no control group and no PC-ablated (ca8$^+$) larvae learned in either the delay or trace paradigms. In contrast, delay ca8$^-$ larvae exhibited robust learning, while trace ca8$^-$ larvae learned to a lesser degree, consistent with the known variability in 3~s trace conditioning performance. These results indicate that intact PCs are required for both delay and trace associative learning, with trace learning showing partial acquisition even in the absence of PCs, reflecting the increased difficulty of the task.

% todo mention increased difficulty of trace conditioning seen in Fig2

% endregion ca8 ablation



% region LTM with puromycin
% ** 3-h Long-Lasting Memory is Not Blocked by Puromycin Under the Current Protocol
\subsection{3-h Long-Lasting Memory is Not Blocked by Puromycin Under the Current Protocol}

To examine the protein synthesis dependence of long-term memory (LTM), I tested delay-conditioned larvae after a 3-h retention interval, with or without the translation inhibitor puromycin (\SI{5}{\milli\gram\per\litre}). Larvae conditioned without puromycin exhibited significant conditioned responses (CRs) during the Early Test phase, showing reduced normalized vigor relative to their Pre-Train baseline and to unpaired controls (Fig.~X). Puromycin-treated larvae displayed a similar reduction in vigor, though this effect did not reach statistical significance, possibly due to the smaller sample size (Fig.~X, top). These results indicate that, under the present experimental conditions, memory expression after 3~h is not blocked by puromycin. This may reflect either a form of long-lasting memory that does not depend on new protein synthesis, or incomplete inhibition of translation at this concentration.

\paragraph{Paradigm Modifications for LTM and Memory Savings Assessment}
To promote LTM formation, we applied spaced training, inserting a \SI{10}{\minute} interval after every 10 training trials (see 5.4. Materials and Methods). A \SI{3}{\hour} stimulus-free retention interval was introduced between training and testing, during which larvae were kept in darkness to minimize sensory stimulation. Baseline tail activity was maintained by a second priming phase immediately before testing, delivering four \SI{100}{\milli\second} US pulses to ensure reliable CR detection.

To test for memory savings, an additional re-training phase (30 paired trials) followed the extinction block (30 CS-only trials). Learning rates during re-training were compared with both the initial acquisition rates and the performance of naive controls trained for the first time.

\paragraph{Three-Hour Memory in Delay Conditioning}
After the 3-h retention interval, delay-trained larvae without puromycin exhibited robust CRs in the Early Test phase, though the response magnitude was smaller than when testing immediately followed training (Fig.~X). Only a subset of individuals (4 of 27) displayed clear CRs at the single-fish level, consistent with partial retention. Puromycin-treated larvae retained comparable CRs, demonstrating that memory at this timescale was not abolished by protein synthesis inhibition. Because dependence on protein synthesis could not be confirmed, I refer to this as \textit{long-lasting memory} rather than canonical LTM.

\paragraph{No Evidence for Memory Savings}
Re-training analysis revealed no significant evidence of memory savings. Previously conditioned larvae reacquired CRs at rates comparable to their initial acquisition and to naive controls trained for the first time (Fig.~X). Naive controls did not reach equally low normalized-vigor values within the first \textasciitilde10 trials, suggesting a possible qualitative difference in early learning dynamics. At the individual level, only a minority of naive controls (4 of 19) developed clear CRs during re-training, indicating that prior extinction did not facilitate faster re-learning under these conditions (Fig.~X).

Together, these results suggest that 3-h post-conditioning memory persists without clear protein synthesis dependence and does not produce detectable savings upon re-exposure to paired stimuli.
% endregion LTM with puromycin

% endregion NMDAR dependence, ca8 ablation, LTM with puromycin





% Imaging figure
%! todo While the main analyses pooled data from fish tested with both white and red CS, a dedicated analysis exclusively on the red CS group confirmed its suitability for subsequent neuroimaging experiments.

